\section{Distribution Reliability Assessment with Failure-Mode Enumeration}
\label{sec:distribution-reliability-fmea}

\subsection{Problem Form and Runtime Scope}

This chapter documents the reliability-assessment formulation for enumerated
contingency analysis on hybrid distribution systems.  The core mathematical
object is a frequency-weighted restoration or service-loss evaluation over a set
of explicit component-failure scenarios rather than pure Monte Carlo sampling.

\subsection{Code Map}

The current implementation and supporting infrastructure are concentrated in:
\begin{itemize}
  \item \texttt{src/analysis/reliability\_assessment.cpp} for reliability-index computation and scenario aggregation,
  \item \texttt{src/analysis/topology\_analysis.cpp} for feeder switching and topology support,
  \item \texttt{src/analysis/distribution\_power\_flow.cpp} for electrical evaluation of restored feeder states.
\end{itemize}

\subsection{Pseudocode Map}

The notebook's detailed algorithm section does not yet carry a dedicated FMEA
pseudocode block.  The closest operational references are the ONR and
distribution-PF workflows in Algorithms A13 and A14 in
Section~\ref{sec:algorithms}.  This chapter therefore serves as the primary
problem-form reference for reliability enumeration itself.

\subsection{Current Status}

This chapter is partly descriptive and partly implementation-facing.  The C++
code base contains the reliability-analysis module, but the documentation here
still plays the role of the canonical formulation for contingency weighting,
duration modeling, and service-restoration metrics.

This section defines a distribution-system reliability framework for hybrid AC/DC
networks that is aligned with the topology-reconfiguration formulation in
\texttt{DistributionPowerFlow/TopologyAnalysis/src/topology\_reconfiguration.jl}
and the current C++ component models.

The target use case is to replace pure Monte Carlo state sampling for
distribution studies with an explicit \(N-1\) failure-mode analysis (FMEA) that:
\begin{enumerate}
  \item enumerates single-component contingencies across AC and DC assets,
  \item models switching and repair windows separately,
  \item co-optimizes restoration dispatch and feeder topology,
  \item quantifies the resilience value of microgrids, network reconfiguration,
        and grid-forming BESS.
\end{enumerate}

\subsection{Sets and Indices}

\begin{align}
\mathcal{N} &= \mathcal{N}^{ac} \cup \mathcal{N}^{dc} && \text{buses}, \\
\mathcal{E} &= \mathcal{E}^{ac} \cup \mathcal{E}^{dc} \cup \mathcal{E}^{vsc} && \text{switchable edges}, \\
\mathcal{G} &= \text{dispatchable generators}, \\
\mathcal{M} &= \text{microgrids (with optional islanding)}, \\
\mathcal{B} &= \text{storage units (including grid-forming BESS)}, \\
\mathcal{L} &= \text{served load points}, \\
\mathcal{C} &= \text{all reliability-tracked components}, \\
\Omega &= \{0\} \cup \{k: k\in\mathcal{C}\} && \text{base state and }N-1\text{ contingencies}, \\
\mathcal{S} &= \{\qual{sw},\qual{rep}\} && \text{switching and repair stages}.
\end{align}

For each contingency \(k\), stage \(s\in\mathcal{S}\), and load \(i\in\mathcal{L}\):
\begin{align}
P^{\qual{shed}}_{i,k,s} &\ge 0 && \text{unsupplied active power (MW)}, \\
E^{\qual{ens}}_{i,k,s} &= P^{\qual{shed}}_{i,k,s}\,\tau_{k,s} && \text{unsupplied energy (MWh)}.
\end{align}

\subsection{Failure and Duration Parameterization}

Each component \(k\in\mathcal{C}\) is represented by annual failure frequency
\(\lambda_k\) (occ/yr) and stage durations
\(\tau_{k,\qual{sw}}\), \(\tau_{k,\qual{rep}}\) (hours).

\paragraph{Field-to-parameter mapping.}
The component models expose either forced outage rate (FOR), failure rate,
or MTBF/MTTR fields. Use the following conversions:
\begin{align}
\lambda_k &= \frac{1}{\text{MTBF}_k}, \\
\mu_k &= \frac{1}{\text{MTTR}_k}, \\
\text{FOR}_k &= \frac{\lambda_k}{\lambda_k + \mu_k}, \\
\lambda_k &= \frac{\text{FOR}_k}{(1-\text{FOR}_k)\,\text{MTTR}_k}, \\
\lambda_k^{\qual{line}} &= \frac{\lambda_k^{\qual{per\,km}}\,\ell_k}{\max(1,n_k^{\qual{parallel}})}.
\end{align}

If a component has scheduled unavailability \(t^{\qual{sched}}_k\) (hr/yr),
include an equivalent availability derating
\(A_k^{\qual{sched}} = 1 - t^{\qual{sched}}_k/8760\), or represent scheduled
outages as deterministic scenarios in \(\Omega\).

\paragraph{Switching duration model.}
For feeder restoration after failure \(k\), define
\begin{align}
\tau_{k,\qual{sw}} = \tau_k^{\qual{detect}} + \tau_k^{\qual{isolate}} + \tau_k^{\qual{reconf}},
\end{align}
where \(\tau_k^{\qual{reconf}}\) may use switch operation data
\(t^{\qual{operation}}\) and switching-failure probability
\(p^{\qual{sw\_fail}}\):
\begin{align}
\tau_k^{\qual{reconf}} = \sum_{r\in\mathcal{R}_k}
\left[(1-p_r^{\qual{sw\_fail}})t_r^{\qual{auto}} + p_r^{\qual{sw\_fail}}t_r^{\qual{manual}}\right].
\end{align}

\subsection{Contingency Enumeration and Weights}

For strict \(N-1\) FMEA, solve one restoration problem per failed component
\(k\in\mathcal{C}\). The expected annual reliability index of any severity term
\(Z_{k,s}\) is frequency-weighted:
\begin{align}
\mathbb{E}[Z] = \sum_{k\in\mathcal{C}} \lambda_k \sum_{s\in\mathcal{S}} Z_{k,s}.
\label{eq:fmea-frequency-weighting}
\end{align}

This avoids Monte Carlo variance and is accurate in the rare-event regime used
in distribution reliability planning.

\subsection{Reconfiguration-Aware Restoration Optimization}

For each contingency \(k\) and stage \(s\), solve an operational restoration
problem with topology and dispatch variables. A linearized mixed-integer model
consistent with the Julia topology formulation is:

\paragraph{Decision variables.}
\begin{align}
\beta_{e,k,s} &\in \{0,1\} && \text{edge energized status}, \\
F_{e,k,s} &\in \mathbb{R} && \text{fictitious/routing flow}, \\
P_{g,k,s},Q_{g,k,s} & && \text{generator dispatch}, \\
P_{m,k,s},Q_{m,k,s} & && \text{microgrid injection}, \\
P_{b,k,s}^{\qual{dis}},P_{b,k,s}^{\qual{ch}},Q_{b,k,s} & && \text{BESS outputs}, \\
P^{\qual{shed}}_{i,k,s} &\ge 0 && \text{load shedding}.
\end{align}

\paragraph{Objective (per contingency-stage).}
\begin{align}
\min\; J_{k,s} =
\sum_{g\in\mathcal{G}} c_g P_{g,k,s}
+ \sum_{b\in\mathcal{B}} c_b^{\qual{cyc}} P_{b,k,s}^{\qual{dis}}
+ \sum_{i\in\mathcal{L}} v_i P^{\qual{shed}}_{i,k,s}
+ c^{\qual{sw}} N^{\qual{sw}}_{k,s}.
\label{eq:restoration-objective}
\end{align}

\paragraph{Topology and radiality constraints.}
\begin{align}
|F_{e,k,s}| &\le M_e\,\beta_{e,k,s}, \quad \forall e\in\mathcal{E}, \\
\sum_{e\in\mathcal{E}} \beta_{e,k,s} &= |\mathcal{N}| - n^{\qual{source}}_{k,s},
\label{eq:radiality-constraint}
\end{align}
with connectivity-flow equations equivalent to the
single-commodity constraints in the network-reconfiguration section.

\paragraph{Failure consistency and reconfiguration logic.}
Let \(\zeta_{e,k}\in\{0,1\}\) denote edge availability under contingency \(k\)
(failed element forced open), and \(\alpha_{e,k,s}\in[0,1]\) the restoration
policy coefficient used in the Julia model:
\begin{align}
\beta_{e,k,s} &\le \zeta_{e,k}, \\
\beta_{e,k,s} &\ge \alpha_{e,k,s} + \zeta_{e,k} - 1.
\label{eq:beta-alpha-zeta}
\end{align}

\paragraph{Nodal power balance with shedding.}
For each bus \(n\in\mathcal{N}\):
\begin{align}
\sum_{g\in\mathcal{G}_n} P_{g,k,s}
+ \sum_{m\in\mathcal{M}_n} P_{m,k,s}
+ \sum_{b\in\mathcal{B}_n} \left(P_{b,k,s}^{\qual{dis}}-P_{b,k,s}^{\qual{ch}}\right)
+ P_{n,k,s}^{\qual{import}}
= P_{n}^{\qual{load}} - P^{\qual{shed}}_{n,k,s}.
\label{eq:nodal-balance-fmea}
\end{align}

Reactive-power and voltage/security constraints follow the same linearized
model family as the ONR and DPF layers.

\subsection{Microgrid and Grid-Forming BESS Formulation}

\subsubsection{Microgrid islanding states}
Define \(z_{m,k,s}\in\{0,1\}\):
\begin{align}
|P_{m,k,s}^{\qual{pcc}}| &\le \overline{P}_m^{\qual{pcc}}(1-z_{m,k,s}), \\
P_{m,k,s}^{\qual{island,gen}} + P_{m,k,s}^{\qual{island,bess}}
&\ge P_{m,k,s}^{\qual{island,load}} - P_{m,k,s}^{\qual{shed}} - M(1-z_{m,k,s}).
\end{align}
When \(z_{m,k,s}=1\), the microgrid supplies internal demand locally;
when \(z_{m,k,s}=0\), PCC exchange is allowed.

\subsubsection{Grid-forming BESS support constraints}
For each BESS \(b\):
\begin{align}
0 \le P_{b,k,s}^{\qual{dis}} &\le \overline{P}_b^{\qual{dis}} u_{b,k,s}, \\
0 \le P_{b,k,s}^{\qual{ch}} &\le \overline{P}_b^{\qual{ch}} (1-u_{b,k,s}), \\
\underline{Q}_b \le Q_{b,k,s} &\le \overline{Q}_b, \\
\text{SOC}_{b,k,s}^{+} &= \text{SOC}_{b,k,s}^{-}
- \frac{\tau_{k,s}}{E_b^{\qual{rated}}}
\left(\frac{P_{b,k,s}^{\qual{dis}}}{\eta_b^{\qual{dis}}}
- \eta_b^{\qual{ch}} P_{b,k,s}^{\qual{ch}}\right), \\
\underline{\text{SOC}}_b \le \text{SOC}_{b,k,s}^{\pm} &\le \overline{\text{SOC}}_b.
\label{eq:bess-soc-outage}
\end{align}

For grid-forming mode, enforce reserve for voltage/frequency support:
\begin{align}
P_{b,k,s}^{\qual{dis}} + R_{b,k,s}^{\qual{up}} \le \overline{P}_b^{\qual{gfm}},
\end{align}
and optionally impose a droop-compatible linear envelope around nominal
frequency/voltage operating points.

\subsection{Reliability Metrics for Distribution Assessment}

\subsubsection{Energy and adequacy indices}
\begin{align}
\text{EENS} &= \sum_{k\in\mathcal{C}} \lambda_k
\sum_{s\in\mathcal{S}} \sum_{i\in\mathcal{L}} E^{\qual{ens}}_{i,k,s}, \\
\text{EDNS} &= \frac{\text{EENS}}{8760}, \\
\text{LOLE} &= \sum_{k\in\mathcal{C}} \lambda_k
\sum_{s\in\mathcal{S}} \tau_{k,s}\,\mathbb{I}\!\left\{\sum_i P^{\qual{shed}}_{i,k,s}>0\right\}, \\
\text{LOLF} &= \sum_{k\in\mathcal{C}} \lambda_k\,
\mathbb{I}\!\left\{\exists s: \sum_i P^{\qual{shed}}_{i,k,s}>0\right\}.
\end{align}

\subsubsection{Customer-based distribution indices}
Let \(N_i\) be customers at load point \(i\), and
\(I_{i,k,s}\in\{0,1\}\) indicate interruption at \((i,k,s)\).
Define nodal interruption frequency and duration:
\begin{align}
\text{CIF}_i &= \sum_{k} \lambda_k \sum_s I_{i,k,s}, \\
\text{CID}_i &= \sum_{k} \lambda_k \sum_s \tau_{k,s} I_{i,k,s}.
\end{align}
Then IEEE 1366 indices are
\begin{align}
\text{SAIFI} &= \frac{\sum_i N_i\,\text{CIF}_i}{\sum_i N_i}, \\
\text{SAIDI} &= \frac{\sum_i N_i\,\text{CID}_i}{\sum_i N_i}, \\
\text{CAIDI} &= \frac{\text{SAIDI}}{\text{SAIFI}}, \\
\text{ASAI} &= 1 - \frac{\text{SAIDI}}{8760}, \\
\text{ASUI} &= 1 - \text{ASAI}.
\end{align}

\subsubsection{Priority-weighted resilience metric}
For critical-service weighting \(w_i\):
\begin{align}
\text{WEENS} = \sum_k \lambda_k \sum_s \sum_i
w_i\,E^{\qual{ens}}_{i,k,s}.
\end{align}

\subsection{Modeling 30+ AC/DC Component Types in One Framework}

The framework uses a unified component-to-contingency mapping:
\begin{align}
\Phi: \mathcal{C} \rightarrow
\left(\lambda_k,\tau_{k,\qual{sw}},\tau_{k,\qual{rep}},\mathcal{U}_k\right),
\end{align}
where \(\mathcal{U}_k\) is the set of model constraints modified by failure \(k\)
(edge outage, generator derating, converter outage, switch lockout, etc.).

The complete field-level catalog is provided in
Section~\ref{sec:distribution-reliability-catalog}.

\subsection{Algorithmic Procedure (Deterministic FMEA)}

\begin{algorithmic}[1]
\State Build component set \(\mathcal{C}\) from AC/DC model tables.
\State Convert reliability fields to \(\lambda_k,\tau_{k,\qual{sw}},\tau_{k,\qual{rep}}\).
\For{each contingency \(k\in\mathcal{C}\)}
  \For{each stage \(s\in\{\qual{sw},\qual{rep}\}\)}
    \State Apply failure state to availability vector \(\zeta_{k}\).
    \State Solve restoration MILP/LP (Eq.\ \ref{eq:restoration-objective}-\ref{eq:nodal-balance-fmea}).
    \State Record \(P^{\qual{shed}}_{i,k,s}\), switching count, and dispatch.
  \EndFor
\EndFor
\State Aggregate EENS/LOLE/LOLF/SAIFI/SAIDI using Eq.\ \ref{eq:fmea-frequency-weighting}.
\State Optionally compare against NSQ/SEQ Monte Carlo as cross-validation.
\end{algorithmic}

\subsection{Implementation Notes for the C++ Reliability Module}

To integrate this formulation into
\texttt{src/analysis/reliability\_assessment.cpp}:
\begin{enumerate}
  \item add a deterministic FMEA execution path alongside NSQ/SEQ Monte Carlo,
  \item add a reconfiguration-aware contingency evaluator (LP/MILP) for
        distribution operation,
  \item include microgrid islanding and BESS SOC update for
        \(\tau_{k,\qual{sw}}\) and \(\tau_{k,\qual{rep}}\),
  \item compute customer-based indices from \texttt{Load::n\_customers} and
        \texttt{Bus::importance}.
\end{enumerate}

This gives a complete bridge from field-level component reliability parameters
to planning-grade reliability metrics under feeder reconfiguration.
